% Transcription — fonds Alexandre Grothendieck, shelfmark 122, batch 2 (pages 21-34).
% Established from the facsimile archives/batches/122/batch-02.pdf (prepared
% from a copy of 122.pdf supplied by the user; PDF page k = archivists' page
% 20+k, no cover sheet), per the skill transcribe-grothendieck. The folder is
% thirty-four pages; this is the second and last batch. Batch 1 (pages 1-20)
% is transcribed separately and did not exist when this pass was made, so the
% notation below is fixed from these pages alone and the triangles (a), (b),
% (c), (d) that these pages invoke are defined on leaves of batch 1.
%
% Pass: Opus 5.5 (claude-opus-5-5), 2026-09-24 — first pass, unchecked against the pages by a human.
%
% Pages skipped: 22, 24, 26, 28, 30, 32 — the versos. They are leaves of an
% English typescript paginated -4.11- to -4.14- (a section 3 on
% counter-examples to geometric quotients: Nagata's quintics, Hironaka's
% example, the blow-ups V_1 ... V_5) and its two hand-drawn figures
% (« Figure 7 », « Figure 8 »), in a hand that is not his. Nothing on them is
% his: skipped by the settled rule, without description in the body.
% No page is summarised.
%
% Pages transcribed: 21, 23, 25, 27, 29, 31 (white paper, blue-black ink) and
% 33, 34 (tan paper). No pagination of his is visible on any leaf.
%
% What the batch is about. Pages 21-31 continue an argument whose beginning is
% in batch 1: W a neighbourhood (a « sphère » in the good case, of dimension
% 2n+1) of an isolated singular point, W_0 a subspace of codimension 2 in it
% (the page's own notation; presumably the special fibre), M ≃ W − W_0 up to
% homotopy, P, Q their pieces, K^• = R(Γ,π)(M) in D(Λ[π]) and L^• =
% RΓ(W_0). Page 21: a triangle (e) for relative cohomology of (W, W_0) →
% (W, Q) → (W, Q̊), a homomorphism of triangles (e) → (b); when W is
% non-singular, Poincaré duality R(Γ,π)_!(M̊) ≃ RHom_Λ(K^•, Λ)[-2n] turns the
% vertical arrow of (a) into u : RHom_Λ(K^•,Λ)[-2n] → K^•, i.e. a class in
% H^{2n}(K^• ⊗^L K^•), and (a) into a triangle (a') whose sides ρ, ρ' are
% transposed (« signe ? »). Page 23: applying the forgetful functor gives
% (c'), applying RΓ^π gives (b'), both « autotransposés »; new data u_1 and
% pairings λ, μ between RΓ(W_0) and RΓ(W − W_0), which in cohomology are the
% maps β and α. Page 25: a pairing RΓ(W_0)^triv × R(Γ,π)(M) → Λ[-(2n-1)];
% (a) and (b) are transposed; an NB that with locally trivial coefficients
% near W_0 the same holds without non-singularity; then W a sphere of
% dimension 2n+1 and the Gysin-type exact sequence. Page 27: the tail of that
% sequence, a boxed criterion for W to be a cohomological sphere, and local
% Lefschetz (X regular at x) with its dual form: the cohomology of W
% determines that of W_0 except in the two middle degrees n-1, n. Page 29:
% vanishing ranges for W a sphere, the short exact pieces, the case n = 1;
% H^i(M, Λ) = 0 for i > n (reading uncertain); a long exact sequence
% H^{2n-i}(M)^∨ → H^i(M) → H^i(W_0) → ... . Page 31: isomorphisms H^i(M) ≃
% H^i(W_0) in low degrees and H^i(W_0) ≃ H^{2n-i-1}(M)^∨ in high degrees;
% for W the topological sphere and n ≥ 3 a boxed two-row exact sequence
% around u in degrees n-1, n. Page 33 (tan): the same statement in the
% language of generic and special fibres V_η, V_s of V — « V sphère
% cohomologique ⟺ » — and, for an isolated singularity, H^i(V_s) = 0 for
% i ≠ 0, n-1, n. Page 34 (tan, another subject): « Inconvénient sérieux » of
% Ext^n(S,T;A) under (AB5) — natural finiteness conditions on a-adic
% inductive systems seem not stable under kernels, and dually under (AB5*) —
% then an « Exemple » (A noetherian, f, g, M, P = M_f/M = lim M/f^n M)
% cancelled by diagonal strokes and broken off at « Q = ».
%
% Notation fixed once for the batch: \mathbb{R}\Gamma, \mathbb{R}\Gamma_{!},
% \mathbb{R}\mathrm{Hom}_\Lambda for his double-struck R; K^\bullet,
% L^\bullet with the bullet he writes as a dot; L^{\bullet\,\mathrm{triv}};
% \mathring{P}, \mathring{Q}, \mathring{M} for his letters with a small ring;
% (-)^\vee for the dual; the circled letters (a), (b), (c), (d), (e), (a'),
% (b'), (c') naming triangles are given as (a) etc., in roman, without circle.
%
% Hardest pages and readings to check first: the diagonal marginal note of
% p. 21 (nearly all \ill{}); the connective prose of pp. 21, 23, 25 (fast
% drafting); the exponent of K^• in (c') (« ens » ?); the end of the top line
% of (b') on p. 23; the overwritten bounds on p. 27 (i ≤ n-2, i = n-1,
% i ≥ n+1, i = n) and on p. 31 (n+2 / n-2, read from « donc 2n-i-1 ≤ n-3 »);
% the boxed bound of p. 29 (i > n or n+1); the subscripts η / ξ on p. 33; the
% whole first paragraph of p. 34.

\documentclass[11pt,a4paper]{article}
% Le préambule commun à toutes les transcriptions du fonds Grothendieck.
%
% Il tient en une idée : **l'incertitude se compose, elle ne se résout pas.**
% Une transcription qui lisse un mot douteux a détruit la seule chose pour
% laquelle on la faisait — un lecteur doit pouvoir savoir, à chaque endroit,
% ce qui était sur le feuillet et ce qui a été ajouté. Les six macros qui
% suivent sont donc l'appareil critique, et non de la décoration.
%
% Les mêmes commandes sont reconnues par `scripts/render.mjs`, qui produit la
% vue de lecture du site. Une macro ajoutée ici doit l'être aussi là-bas,
% faute de quoi le rendu échouera bruyamment — ce qui est voulu : un rendu
% silencieusement faux est pire qu'un rendu absent.

\ProvidesPackage{grothendieck}[2026/08/08 Transcriptions du fonds Grothendieck]

\usepackage{amsmath,amssymb,amsthm}
\usepackage{mathrsfs}
\usepackage{stmaryrd}
% Les diagrammes commutatifs sont partout dans ce fonds : c'est la langue
% même de la géométrie algébrique de Grothendieck, et une transcription qui
% les rendrait en prose ne transcrirait rien.
\usepackage{tikz-cd}
% Français par défaut — c'est la langue des feuillets — mais l'anglais est
% chargé aussi : la traduction et le résumé sont des documents anglais, et la
% typographie française (espaces avant les deux-points) y serait une faute.
% Ces documents basculent par \selectlanguage{english} en tête de corps.
\usepackage[english,french]{babel}
\usepackage{fontspec}
\usepackage{geometry}
\geometry{a4paper,margin=25mm,marginparwidth=28mm}
\usepackage{marginnote}
\usepackage{xcolor}
% ulem, not soul. `soul`'s \st and \ul are built on a character-by-character
% scan that XeLaTeX's font machinery breaks: the document fails with an
% undefined control sequence rather than degrading. ulem does the same two jobs
% and is Unicode-engine safe. `normalem` stops it hijacking \emph, which the
% transcriptions use for Grothendieck's own emphasis.
\usepackage[normalem]{ulem}
\usepackage{hyperref}
\hypersetup{colorlinks=true,linkcolor=black,urlcolor=black,pdfborder={0 0 0}}

% --- Métadonnées du lot -----------------------------------------------------
% Reprises telles quelles par le script de rendu, qui les affiche en tête de
% la vue de lecture. Elles fixent ce que le fichier prétend couvrir : sans
% elles, une transcription flotte sans point d'attache dans un fonds de seize
% mille feuillets.

\newcommand{\folder}[1]{\gdef\grothFolder{#1}}
\newcommand{\batch}[1]{\gdef\grothBatch{#1}}
\newcommand{\leaves}[2]{\gdef\grothFirst{#1}\gdef\grothLast{#2}}
\newcommand{\foldertitle}[1]{\gdef\grothFoldertitle{#1}}
\gdef\grothFolder{?}\gdef\grothBatch{?}\gdef\grothFirst{?}\gdef\grothLast{?}\gdef\grothFoldertitle{}

% --- Le résumé --------------------------------------------------------------
%
% Il ouvre la lecture modernisée, sous le titre « Résumé », et il est composé
% en petit. Ce n'est pas un ornement : c'est le seul passage du document qui
% ne soit pas des mathématiques, et un lecteur doit pouvoir le reconnaître
% avant de l'avoir lu — puis le sauter s'il connaît déjà le sujet, ou s'y
% arrêter s'il ne le connaît pas. Le filet qui le suit marque la reprise du
% corps.

\newenvironment{resume}
  {\small\setlength{\parskip}{0.45em}}
  {\par\medskip\hrule\medskip}

% --- Mots-clés ---------------------------------------------------------------
%
% En anglais, en clôture du résumé de la lecture modernisée : le vocabulaire
% moderne sous lequel chercher ce que les feuillets construisent. Le manifeste
% du site (`npm run manifest`) les extrait pour étiqueter la cote entière —
% cette ligne est la source unique des tags, il n'y a pas de fichier à côté.
\newcommand{\keywords}[1]{\par\smallskip\noindent
  {\footnotesize\textsc{Keywords} --- \textit{#1}}\par}

% --- L'appareil critique ----------------------------------------------------

% Le numéro de feuillet, en marge. C'est l'ancre qui permet de régler un
% désaccord sur une formule en regardant *un* feuillet plutôt qu'une chemise
% de six cents. Le site s'en sert aussi pour tourner les pages du fac-similé
% quand on fait défiler la transcription.
\newcommand{\leaf}[1]{%
  \marginnote{\footnotesize\color{gray!70}\textsf{#1}}%
  \label{leaf:#1}\ignorespaces}

% Illisible : on ne devine pas. Un mot inventé qui se lit comme les autres est
% le pire résultat possible.
\newcommand{\ill}{\textcolor{red!60!black}{[\,\ldots\,]}}

% Lecture douteuse : le mot est donné, et signalé comme incertain.
\newcommand{\uncertain}[1]{\uline{#1}}

% Ajout de l'éditeur — une lettre manquante, un mot sous-entendu.
\newcommand{\add}[1]{\textcolor{blue!50!black}{[#1]}}

% Biffé par Grothendieck lui-même. Ce qu'il a rayé fait partie du document :
% une rature dit souvent où la pensée a changé de direction.
\newcommand{\struck}[1]{\sout{#1}}

% Note du transcripteur, sur l'état matériel du feuillet.
\newcommand{\note}[1]{\par\smallskip{\small\color{gray!60!black}\textit{[#1]}}\par\smallskip}

% Note marginale de Grothendieck, distincte du corps du texte.
\newcommand{\marginal}[1]{\par\smallskip\noindent\hspace{1em}%
  {\small\color{gray!60!black}#1}\par\smallskip}

% --- Titre ------------------------------------------------------------------

\AtBeginDocument{%
  \begin{center}
    {\large\bfseries Fonds Alexandre Grothendieck --- cote n\textsuperscript{o}~\grothFolder}\\[2pt]
    {\itshape\grothFoldertitle}\\[4pt]
    {\small feuillets \grothFirst--\grothLast\ (lot \grothBatch)}\\[2pt]
    {\footnotesize\color{gray!60!black}Transcription établie d'après le fac-similé numérisé
     par l'Université de Montpellier.\\
     Les lectures incertaines sont soulignées, les illisibles notées [\,\ldots\,],
     les ajouts entre crochets.}
  \end{center}
  \vspace{1em}}

\endinput


\folder{122}
\batch{2}
\pages{21}{34}
\dating{[vers 1968-1971]}
\watermark{Édition de démonstration}
\foldertitle{Point singulier isolé de la fibre : notes manuscrites (s.d.)}

\begin{document}

\page{21}

(e)

\begin{tikzcd}[column sep=small]
 & \mathbb{R}\Gamma(W_0) \simeq \mathbb{R}\Gamma(Q) \simeq \mathbb{R}\Gamma(\mathring{Q}) \arrow[dl, dashed, "\beta'"'] \\
\mathbb{R}\Gamma_{!}(\mathring{P}) \arrow[r, dashed] & \mathbb{R}\Gamma(W) \arrow[u, "\rho_1"']
\end{tikzcd}

\note{les flèches $\beta'$ et $\mathbb{R}\Gamma_{!}(\mathring{P}) \to
\mathbb{R}\Gamma(W)$ sont tracées en pointillé ; le sommet
$\mathbb{R}\Gamma_{!}(\mathring{P})$ est mis entre crochets et entouré au crayon.
Le long de $\rho_1$ : « \uncertain{compatible avec} \ill{} \ill{} ». À droite :
« NB $\rho_1 = \alpha\rho_0$ ».}

[suite exacte de cohomologie relative pour $(W, W_0) \to (W, Q) \to (W, \mathring{Q})$].

\textbf{NB}. On a un homom\add{orphisme} $(u_0, u_1, u_2)$ des triangles (e) dans les
triangles (b), \struck{\ill{}} \uncertain{induit} par $(P, \partial P) \subset (W, Q)$,
avec \note{sous le $Q$, en interligne : « $W_0$ », suivi d'un mot biffé}
\[
  u_1 = \rho_0, \qquad u_2 = \mathrm{id}_1
\]
et qui induit un isom\add{orphisme} $u_0$ (grâce aux \ill{} d'identification
\ill{} \ill{} \ill{} $=$ \ill{} \ldots).

\marginal{\ill{} \ill{} le cas \ill{} \ill{} \ill{} \ill{} \ill{} $W$ \ill{}
\ill{} \ill{} (a') \ill{}}\note{note écrite en oblique dans la marge gauche, en
regard des lignes précédentes ; presque entièrement illisible}

Supposons $W$ non singulière, i.e. $\mathring{M}$ non singulière, i.e.
$\mathring{P}$ non singulière. Alors on a un isomorphisme
\[
  \mathbb{R}(\Gamma, \pi)_{!}(\mathring{M}) \simeq
  \mathbb{R}\mathrm{Hom}_{\Lambda}(K^\bullet, \Lambda)[2n],
  \qquad K^\bullet = \mathbb{R}(\Gamma, \pi)(M),
\]
\note{le décalage est écrit $[2n]$, le crochet ouvrant repassé ; plus bas
sur la page il est $[-2n]$. L'égalité $K^\bullet = \mathbb{R}(\Gamma,\pi)(M)$
est écrite sous $K^\bullet$}
et la flèche verticale en (a) est une flèche
\[
  \boxed{\ u : \mathbb{R}\mathrm{Hom}_{\Lambda}(K^\bullet, \Lambda)[-2n] \longrightarrow K^\bullet\ }
  \qquad \text{dans } D(\Lambda[\pi]),
\]
i.e. un élément de \struck{\ill{}}
\[
  \xi \in H^{2n}\bigl(K^\bullet \overset{\mathbf{L}}{\otimes}_{\Lambda} K^\bullet\bigr).
\]
\note{la lettre $\xi$ sous l'élément est douteuse}

La donnée \uncertain{précédente} (a) devient

(a')

\begin{tikzcd}[column sep=small]
 & L^{\bullet\,\mathrm{triv}} \arrow[dl, "\rho'"'] & \\
\mathbb{R}\mathrm{Hom}_{\Lambda}(K^\bullet, \Lambda)[-2n] \arrow[rr] & & K^\bullet \arrow[ul, "\rho"']
\end{tikzcd}

\marginal{$K^\bullet \in \mathrm{ob}\, D(\Lambda[\pi])$, $L^\bullet \in \mathrm{ob}\, D(\Lambda)$,
$K^\bullet$ parfait rel\add{ativement à} $\Lambda$}

En fait, $\rho$ et $\rho'$ sont \emph{transposés} l'un de l'autre (signe ?)
dans les triangles (a) et (a') \ill{} comme \uncertain{autotransposé}.

\page{23}

Appliquant le foncteur oubli, on trouve

(c')

\begin{tikzcd}[column sep=small]
 & L^\bullet \simeq \mathbb{R}\Gamma(W_0) \arrow[dl] & \\
\mathbb{R}\mathrm{Hom}_{\Lambda}(K^{\bullet\,\mathrm{ens}}, \Lambda)[-2n] \arrow[rr] & & K^{\bullet\,\mathrm{ens}} \simeq \mathbb{R}\Gamma(M) \arrow[ul]
\end{tikzcd}

\note{l'exposant de $K^\bullet$, lu « ens », est douteux. Sous le premier
sommet : $\mathbb{R}\mathrm{Hom}_{\Lambda}(K^{\bullet\,\mathrm{ens}}, \Lambda)[-2n]
\simeq \mathbb{R}\Gamma_{!}(\mathring{M})$}

\marginal{dire que c'est encore un triangle autotransposé}

Appliquant $\mathbb{R}\Gamma^{\pi}$, on trouve [\uncertain{exprimer} comment ?]

(b')

\begin{tikzcd}[column sep=small, row sep=small, nodes={font=\scriptsize}]
 & L^\bullet \overset{\mathbf{L}}{\otimes}_{\Lambda} \mathbb{R}\Gamma(S^1) \simeq L^\bullet \times L^\bullet[-1] \arrow[dl, "{(\alpha',\ \beta')}"'] & \\
\mathbb{R}\mathrm{Hom}_{\Lambda}(K^{\bullet\,\pi}, \Lambda)[-2n-1] \arrow[rr] & & K^{\bullet\,\pi} \simeq \mathbb{R}\Gamma(P) \simeq \mathbb{R}\Gamma(W - W_0) \arrow[ul, "{(\alpha,\ \beta)}"']
\end{tikzcd}

\note{le sommet supérieur se prolonge par « $\simeq$ \ill{} ». Sous le sommet
de gauche : $\mathbb{R}\Gamma_{!}(\mathring{P}) \simeq \mathbb{R}\Gamma_{!}(W - W_0)$}

(triangle autotransposé)

\marginal{vérifier la commutation de $\mathbb{R}\Gamma^{\pi}$ à
$\mathbb{R}\mathrm{Hom}_{\Lambda}(-, \Lambda) = (-)^\vee$}\note{dans un
cadre, à gauche, relié au triangle (b')}

\struck{\ill{} les triangles (d) et (e) \ill{} \ill{}}\note{ligne
entièrement biffée}

Comme nouvelles \struck{données} \uncertain{structures}, on trouve \ill{} deux
\struck{$u_1 : \mathbb{R}\mathrm{Hom}_{\Lambda}(K^\bullet, \Lambda)[-2n-1] \longrightarrow$}

\[
  \boxed{\ u_1 : \mathbb{R}\Gamma(W - W_0)^\vee[-2n-1] \longrightarrow \mathbb{R}\Gamma(W - W_0)\ }
\]
et des accouplements
\[
  \boxed{\begin{aligned}
    \lambda &: \mathbb{R}\Gamma(W_0) \times \mathbb{R}\Gamma(W - W_0) \longrightarrow \Lambda[-2n] \\
    \mu &: \mathbb{R}\Gamma(W_0) \times \mathbb{R}\Gamma(W - W_0) \longrightarrow \Lambda[-(2n-1)]
  \end{aligned}}
\]
\[
  \begin{aligned}
    \lambda &: H^{2n-i}(W_0) \times H^{i}(W - W_0) \longrightarrow \Lambda \\
    \mu &: H^{2n-1-i}(W_0) \times H^{i}(W - W_0) \longrightarrow \Lambda
  \end{aligned}
\]
\note{les exposants de ces deux lignes sont récrits par-dessus d'autres, illisibles}
i.e.
\[
  \begin{aligned}
    \lambda^i &: H^{i}(W - W_0) \longrightarrow H^{i-1}(W_0) && \text{je dis que c'est } \beta \\
    \mu^i &: H^{i}(W - W_0) \longrightarrow H^{i}(W_0) && \text{je dis que c'est } \alpha
  \end{aligned}
\]

\page{25}

On obtient donc deux données remarquables, savoir $u$ et
\[
  L^{\bullet\,\mathrm{triv}} \times K^\bullet \longrightarrow \Lambda[-2n+1]
\]
i.e.
\[
  \boxed{\ \text{accoupl\add{ement}}\quad \mathbb{R}\Gamma(W_0)^{\mathrm{triv}} \times
  \mathbb{R}(\Gamma, \pi)(M) \longrightarrow \Lambda[-(2n-1)]\ }
\]
i.e.
\[
  \mathbb{R}(\Gamma, \pi)(M) \longrightarrow
  \mathbb{R}\Gamma(W_0)^\vee[-(2n-1)]^{\mathrm{triv}} \simeq \mathbb{R}\Gamma(W_0)^{\mathrm{triv}}
\]
\note{sous le dual : « p.r. à $\Lambda$ » ; sous l'isomorphisme : « dualité de $W_0$ »}

\emph{Enfin}, les \emph{triangles} (a) \emph{et} (b) \emph{sont transposés}.

\textbf{NB}. En prenant des coefficients généraux, loc\add{alement}
triviaux au voisinage de \struck{\ill{}} $W_0$, on trouve que, prenant des
coefficients \ill{} : des coefficients « duaux », \struck{les \ill{}} sans
hyp\add{othèse} de \struck{régularité} non-singularité sur $W$, les suites
exactes du type (a), (b), (e) sont transposées de suites exactes du même type,
et les suites exactes du type (c) sont transposées de celles du type (d).
\note{la troisième lettre de la première liste est surchargée ; (e) est
une lecture douteuse}

\marginal{plus bas}

Supposons que $W$ soit \emph{une sphère} (de dim\add{ension} $2n+1$). Alors
la suite exacte \struck{de} \uncertain{déduite de} (b)
\[
  \cdots \longrightarrow H^{i-2}(W_0) \longrightarrow H^{i}(W) \longrightarrow
  H^{i}(W - W_0) \longrightarrow H^{i-1}(W_0) \longrightarrow H^{i+1}(W)
\]
\uncertain{montre ceci} : i.e. $i \neq 0, 2n, 2n+1$ (en fait, comme on sait
(b), il suffit que $i \neq 0, 2n$)

a) si $i, i+1 \neq 0, 2n+1$, alors $H^{i}(W - W_0) \xrightarrow{\ \sim\ } H^{i-1}(W_0)$

\struck{b) $\to H^{2n}(W - W_0) \to H^{2n-1}(W_0) \to \Lambda \to \ldots \to
H^{2n}(W_0) \to 0$}

\struck{$H^{2n+1}(W - W_0) \simeq H^{2n}(W_0) \simeq 0$}

\page{27}

b) La portion de suite exacte
\[
  \begin{aligned}
  &\underbrace{H^{2n}(W)}_{=0} \to H^{2n}(W - W_0) \xrightarrow{\ \mathrm{inj}\ }
  \underbrace{H^{2n-1}(W_0)}_{\simeq \Lambda} \xrightarrow{\ \mathrm{surj}\ }
  \underbrace{H^{2n+1}(W)}_{=\Lambda} \\
  &\qquad \to \underbrace{H^{2n+1}(W - W_0)}_{=0}
  \to \underbrace{H^{2n}(W_0)}_{=0} \to 0
  \end{aligned}
\]
\note{sous $H^{2n-1}(W_0) \simeq \Lambda$, dans une bulle : « si $W_0$ est
connexe, \ill{} le cas si $n \geq 2$ »}
donne :
\[
  \left\{\begin{aligned}
    H^{2n}(W - W_0) &= H^{2n+1}(W - W_0) = 0 \\
    \Lambda = H^{2n-1}(W_0) &\simeq H^{2n+1}(W) = \Lambda
  \end{aligned}\right.
\]
\marginal{comme on \ill{} \uncertain{résulte} de Lefschetz \uncertain{affine?} local plus bas}\note{dans
une bulle reliée à la première égalité}

\marginal{NB $H^{i}(W_0) = 0$ sauf si $i \in [0, 2n-1]$. $H^{i}(W_0)$ satisfait
dualité de Poincaré}\note{encadré à droite}

\marginal{NB $W$ sph\add{ère} coh\add{omologique} $\Longleftrightarrow$
$\bigl\{H^{i}(W - W_0) \xrightarrow{\sim} H^{i-1}(W_0)$ si $i \neq 0, 2n$
\emph{et} $H^{2n}(W - W_0) \hookrightarrow H^{2n-1}(W_0)$ [si \struck{\ill{}}
$W_0$ connexe, cela équivaut à $H^{2n}(W - W_0) = 0$]. Si $n \geq 2$, et $W$
non singulière, il \uncertain{faut} que $H^{2n}(W - W_0) = 0$, $W_0$
connexe$\bigr\}$}\note{encadré à gauche, sous un mot fortement biffé}

Supposons que $X$ soit régulier en $x$. Alors on a, par
\emph{Lefschetz local} :
\[
  \left\{\begin{aligned}
    H^{i}(W) \longrightarrow H^{i}(W_0) \quad &\text{un isom pour } i \leq n-2 \\
    &\text{un mono pour } i = n-1
  \end{aligned}\right.
\]
\note{les deux bornes sont récrites par-dessus d'autres chiffres}

\struck{Dire que $W$ est une sphère,}

Par \struck{la suite exacte \uncertain{cohomologie} si et} dualité, cela signifie
\[
  \left\{\begin{aligned}
    H^{i}_{!}(W_0) \longrightarrow H^{i+2}(W) \quad &\text{un isom si } i \geq n+1 \\
    &\text{un épim si } i = n
  \end{aligned}\right.
\]
\note{les bornes $n+1$ et $n$ sont surchargées ; la seconde est douteuse}
donc la connaissance de la cohomologie de $W$ implique celle de la
cohomologie de $W_0$, sauf en les deux dimensions \struck{moyennes}
$H^{n-1}(W_0)$, $H^{n}(W_0)$ (complémentaires, i.e. où les cohomologies sont
duales pour $W_0$).

On peut dire aussi que l'on a
\[
  \Bigl[\ H^{i}(W - W_0) = 0 \ \text{ si } i > n+1 \Bigr],
  \qquad H^{i}(W - W_0) = H^{i}(P).
\]
\note{l'égalité $H^{i}(W - W_0) = H^{i}(P)$ est écrite verticalement sous le
premier terme}

\page{29}

Revenant au cas $W$ une sphère, on trouve
\[
  \left\{\begin{aligned}
    H^{i}(W - W_0) &= 0 \quad \text{si } i > n+1 \\
    H^{i}(W_0) &= 0 \quad \text{si } i \neq n-1, n, 0, 2n-1
  \end{aligned}\right.
\]
[en fait, si $i \neq$ \struck{\ill{}} $0, 1, n, n+1$] (d'abord si $i \geq n+1$
et $i \neq 2n-1$, puis le reste par dualité)

La suite exacte \ill{} : \ill{} est triviale, sauf la portion
\[
  0 \to H^{0}(W) \to H^{0}(W - W_0) \to 0, \qquad
  0 \to H^{1}(W - W_0) \to H^{0}(W_0) \to 0
\]
\[
  \left\{\begin{aligned}
    &0 \to H^{n}(W - W_0) \to H^{n-1}(W_0) \to 0 \\
    &0 \to H^{n+1}(W - W_0) \to H^{n}(W_0) \to 0 \qquad (\text{si } n \neq 1) \\
    &0 \to H^{2n-1}(W_0) \to H^{2n+1}(W) \to 0
  \end{aligned}\right.
\]
\note{à droite, deux lignes biffées : « \struck{$n+1 \neq 2n+1$} »,
« \struck{$n+2 \neq 2n+1$} » (lecture douteuse)}

[Si $n = 1$, on trouve
\[
  \begin{aligned}
    &0 \to H^{0}(W) \xrightarrow{\ \sim\ } H^{0}(W - W_0) \to 0 \\
    &0 \to H^{1}(W - W_0) \to H^{0}(W_0) \to 0 \\
    &0 \to H^{2}(W - W_0) \to H^{1}(W_0) \to H^{3}(W) \simeq \Lambda \to 0 \ ]
  \end{aligned}
\]

Notons que, pour \ill{} : la \ill{}, \struck{\ill{}} les relations
$H^{i}(W - W_0, \Lambda) = 0$ pour $i > n+1$ donnent
\[
  \boxed{\ H^{i}(M, \Lambda) = 0 \quad \text{si } i > n\ }
\]
\note{la borne de l'encadré est surchargée ; $n$ est douteux, $n+1$ possible.
Dans la ligne précédente, $W - W_0$ est récrit sur un autre symbole}
[Indépendant en fait \ill{} que $W$ non singulière et de ce que $W_0$ non singulière]

La suite exacte infinie (si $W$ non singulière)
\[
  \cdots \to H^{2n-i}(M)^\vee \to H^{i}(M) \to H^{i}(W_0) \to
  H^{2n-i-1}(M)^\vee \to H^{i+1}(M) \to H^{i+1}(W_0)
\]
montre donc que \struck{si $i > n+1$}

\page{31}

\struck{En d'autres termes, le triangle \ldots\ $H^{i}(W_0) = 0$ si $i \neq$ \ill{}}\note{deux
lignes biffées, avec un petit encadré « $i \neq n-1$ », « $2n-1-i = n-1$ »}

\marginal{isomorphismes}
\[
  \left\{\begin{aligned}
    H^{i}(W_0) &\xrightarrow{\ \sim\ } H^{2n-i-1}(M)^\vee && \text{si } i \geq n+2
      && (\text{épim si } i = n+1) \quad (\text{donc } 2n-i-1 \leq n-3) \\
    H^{i}(M) &\xrightarrow{\ \sim\ } H^{i}(W_0) && \text{si } i \leq n-2
      && (\text{mono si } i = n-1)
  \end{aligned}\right.
\]
\note{les deux bornes sont surchargées d'un chiffre épais ; $n+2$ et $n-2$
sont lus d'après « donc $2n-i-1 \leq n-3$ » et la page~27. La seconde ligne
est entourée. Les mots « épim » et « mono » entre parenthèses sont des lectures douteuses}

Donc $H^{q}(M)$ est connu en termes de $H^{q}(W_0)$ sauf en dimensions
\struck{\ill{}} $n-1$, $n$, \struck{\ill{}}, où on a la suite exacte

\begin{tikzcd}[column sep=small]
H^{n-1}(M) \arrow[r] & H^{n-1}(W_0) \arrow[r, "\mathrm{surj}"] & H^{n}(M)^\vee \arrow[d] \\
H^{n-1}(M)^\vee & H^{n}(W_0) \arrow[l] & H^{n}(M) \arrow[l]
\end{tikzcd}

\note{les deux lignes commencent par des termes biffés :
$\to H^{n-2}(M) \to H^{n-2}(W_0) \to H^{n+1}(M)^\vee\ (=0) \to$ en haut,
$0 \leftarrow H^{n-2}(M)^\vee \leftarrow H^{n-1}(W_0) \leftarrow H^{n-1}(M)\ (=0) \leftarrow$
en bas ; entre les deux lignes, un signe biffé}

Supposons maintenant que $W$ \ill{} la sphère topologique. Alors on a vu que
$H^{i}(W_0) = 0$ si $i \neq 0, n-1, n, 2n-1$.

\emph{Si donc} $n \geq 3$, on trouve aussi \emph{$H^{i}(M) = 0$},
\struck{\ill{}}\note{en interligne : « pour $i \neq 0, n-1, n$,
\struck{\ill{}} »}

\struck{et la suite exacte précédente devient}
\[
  \boxed{\begin{array}{c}
    0 \to H^{n+1}(M)^\vee \xrightarrow{\ u\ } H^{n-1}(M) \to H^{n-1}(W_0) \to H^{n}(M)^\vee \\
    \downarrow \\
    0 \leftarrow H^{n+1}(M) \xleftarrow{\ {}^t u\ } H^{n-1}(M)^\vee \leftarrow H^{n}(W_0) \leftarrow H^{n}(M)
  \end{array}}
\]
\note{la flèche verticale descend de $H^{n}(M)^\vee$ à $H^{n}(M)$, en bout de
ligne. L'encadré est traversé au crayon de traits obliques en zigzag reliant
les deux lignes, avec des signes $+$ et $-$ alternés au-dessus et au-dessous
des termes ; l'exposant du premier terme de la ligne supérieure est surchargé}

\page{33}

$V$ une sphère cohomologique
\[
  \Updownarrow
\]
\[
  \left\{\begin{aligned}
    &H^{i}(V_\eta) \to H^{i-1}(V_s) \ \text{un \emph{isom} si } i \neq 0, 2n \\
    &\emph{et}\ \left\{\begin{aligned}
      &H^{0}(V_\eta) \simeq \Lambda \\
      &H^{2n}(V_\eta) \to H^{2n-1}(V_s) \ \text{inj et conoyau} \simeq \Lambda \\
    \end{aligned}\right. \qquad (\text{quand } n \geq 1)
  \end{aligned}\right.
\]
\note{avant « $i \neq 0, 2n$ », biffé : « \struck{$i, i+1 \neq 0, 2n+1$} » ; après, un terme biffé illisible ; la troisième ligne de l'accolade intérieure est biffée : « \struck{$H^{2n+1}(V_s) \simeq \Lambda$} ». Sous la ligne : « $2n+1 > n+1$ »}
\[
  \Bigl[\ 0 \to \underbrace{H^{2n}(V_\eta)}_{=0 \text{ si } n \geq 2} \to H^{2n-1}(V_s)
  \to H^{2n+1}(V) \to \underbrace{H^{2n+1}(V_\eta)}_{=0 \text{ si } n \geq 1}
  \to \underbrace{H^{2n}(V_s)}_{=0} \to 0 \ \Bigr]
\]
\note{les exposants de $H^{2n+1}(V)$ et $H^{2n+1}(V_\eta)$ se lisent aussi
bien $2n-1$}
\[
  \Downarrow \quad \text{si on suppose déjà que l'on a singularité « isolée »}
\]
\[
  H^{i}(V_s) = 0 \quad \text{si } i \neq 0, n-1, n
\]
\marginal{$H^{i}(V_{\xi}) \xrightarrow{\ \sim\ } H^{i}(V_s)$}\note{à droite,
séparé par un trait vertical ; l'indice, lu $\xi$, peut être le même que
l'$\eta$ plus haut}

\page{34}

\emph{Inconvénient sérieux} de la \ill{}
\uncertain{Ext}$^n(S, T; A)$ lorsque $A$ satisfait (AB5), \struck{\ill{}}
\ill{} : les conditions de finitude naturelles (pour $A$ loc\add{alement}
noeth\add{érien}), sur les systèmes inductifs $a$-adiques, ou ess\add{entiellement}
$a$-adiques, ou \ill{} en $a$-adiques, ne semblent pas stables par
\emph{\uncertain{noyaux}} [dualement, des \ill{}$(S, T; A)$ quand $A$
satisfait (AB5$^*$), les conditions de finitude \ill{} ne semblent pas
stables par \uncertain{conoyaux}] !!

\struck{\emph{Exemple}. $A$ anneau [noeth.], $f, g$ deux éléments de $A$,
$M$ un objet de \ill{} $\mathcal{C}$, \ill{} $(f, g)$ soit $M$-régulière. Soit
$P = M_f / M = \varinjlim_n M / f^n M$ (i.e. $M/f^nM \to M/f^mM$ est donné par
$f^{m-n}$ ; le système inductif est à morphismes de transition injectifs.)
Soit $Q =$}\note{le paragraphe est barré de traits obliques croisés et
s'arrête sur « $Q =$ »}

\end{document}
